\documentclass[conference]{IEEEtran}
\IEEEoverridecommandlockouts

\usepackage{cite}
\usepackage{amsmath,amssymb,amsfonts}
\usepackage{algorithmic}
\usepackage{graphicx}
\usepackage{textcomp}
\usepackage{xcolor}
\usepackage{booktabs}
\usepackage{multirow}
\usepackage{hyperref}
\usepackage{balance}

\def\BibTeX{{\rm B\kern-.05em{\sc i\kern-.025em b}\kern-.08em
    T\kern-.1667em\lower.7ex\hbox{E}\kern-.125emX}}

\begin{document}

\title{When AI Re-enacts Science: A Meta-Study on Large Language Model Agents Reproducing Machine Learning Research\\
{\footnotesize \textsuperscript{*}Implications for Engineering Research Reproducibility}\\[0.5em]
{\small Version 2.0 -- Major Revision}}

\author{\IEEEauthorblockN{Claude Code\textsuperscript{1} and Fredrik Ahlgren\textsuperscript{2}}
\IEEEauthorblockA{\textsuperscript{1}Anthropic, San Francisco, CA, USA (AI Research Agent)\\
\textsuperscript{2}Human Supervisor, Sweden}}

\maketitle

\begin{abstract}
The reproducibility crisis in scientific research raises fundamental questions about verification and validation of published findings. This paper presents a meta-study examining what happens when a large language model (LLM) agent attempts to \textit{re-enact} engineering research from documentation alone. Using Claude Code to re-enact a 2018 study on ship fuel consumption prediction---with synthetic data generated from documented statistics since original data was unavailable---we investigate what the process reveals about documentation sufficiency, methodological choices, and AI-driven verification. The re-enactment achieved R\textsuperscript{2} = 0.9921 $\pm$ 0.0008 on synthetic data (vs. original 0.992), while polynomial feature engineering achieved R\textsuperscript{2} = 0.9963 $\pm$ 0.0005, suggesting unexplored alternatives. Critically, we distinguish this \textit{documentation-driven re-enactment} from true replication on original data. Comparison with recent LLM reproduction benchmarks (SciReplicate-Bench: 39\% success) suggests task characteristics---interpretable physics, available code, clear statistics---enabled our success. We propose an ``AI replicability checklist'' for research documentation and discuss implications for verification at scale.
\end{abstract}

\begin{IEEEkeywords}
scientific reproducibility, artificial intelligence, machine learning, AutoML, meta-research, engineering methodology, research verification
\end{IEEEkeywords}

\section{Introduction}

\subsection{The Reproducibility Crisis}

Scientific reproducibility---the ability to obtain consistent results using the same data and methods---is foundational to the scientific enterprise \cite{baker2016reproducibility}. Yet surveys indicate that over 70\% of researchers have failed to reproduce another scientist's experiments, and more than half have failed to reproduce their own \cite{baker2016reproducibility}. In machine learning specifically, Hutson \cite{hutson2018artificial} documented systematic failures in reproducing published results, attributing this to incomplete documentation of hyperparameters, data preprocessing, and computational environments.

The implications extend beyond academic integrity. Engineering research informs real-world systems---from autonomous vehicles to energy infrastructure. If published methods cannot be reliably reproduced, the foundation for technological progress becomes uncertain.

\subsection{Terminology: Reproducibility, Replicability, Re-enactment}

Following the ACM and recent reproducibility literature \cite{barba2018terminologies}, we distinguish:

\begin{itemize}
    \item \textbf{Reproducibility}: Same data, same code $\rightarrow$ same results
    \item \textbf{Replicability}: New data, same methods $\rightarrow$ consistent results
    \item \textbf{Re-enactment}: Synthetic data from documentation $\rightarrow$ tests documentation sufficiency
\end{itemize}

This study performs \textit{documentation-driven re-enactment}: we reconstruct experiments from published descriptions using synthetic data that matches documented statistics. This tests whether documentation is sufficient to enable independent verification---a distinct but valuable question from whether original empirical claims are reproducible.

\subsection{Research Questions}

This meta-study addresses: \textbf{What can AI-driven re-enactment reveal about research documentation and methodology?}

We decompose this into specific inquiries:

\begin{enumerate}
    \item Can an LLM agent successfully re-enact ML research methodology from documentation alone?
    \item What does the process reveal about documentation quality and implicit assumptions?
    \item How do results contextualize within modern methodological advances?
    \item What are the implications for research verification at scale?
\end{enumerate}

\subsection{Contributions}

This paper makes four primary contributions:

\begin{enumerate}
    \item \textbf{Documentation-Driven Re-enactment}: We demonstrate that an LLM agent can re-enact ML research from documentation, achieving R\textsuperscript{2} = 0.9921 $\pm$ 0.0008 on synthetic data matching original statistics.

    \item \textbf{Critical Analysis Framework}: We show that AI re-enactment provides systematic critique---testing 14 alternative methods, identifying documentation gaps, and contextualizing 2018 choices within the modern landscape.

    \item \textbf{Self-Replication Perspective}: Uniquely, the human supervisor (Ahlgren) is the original paper's first author, providing insider perspective on documentation sufficiency and implicit assumptions.

    \item \textbf{AI Replicability Checklist}: We propose concrete documentation standards based on gaps identified during re-enactment.
\end{enumerate}

\section{Background and Related Work}

\subsection{The Original Study}

The target paper for replication is ``Auto Machine Learning for predicting Ship Fuel Consumption'' by Ahlgren and Thern \cite{ahlgren2018automl}, presented at ECOS 2018. The study applied TPOT \cite{olson2016tpot}, a genetic algorithm-based AutoML tool, to predict fuel oil consumption on a Baltic Sea cruise ship using engine sensor data.

Key characteristics of the original work:
\begin{itemize}
    \item \textbf{Data}: 10 months of sensor data from 8 engines (4 auxiliary, 4 main)
    \item \textbf{Features}: Engine RPM, fuel rack position, exhaust temperature, turbocharger RPM
    \item \textbf{Target}: Fuel oil consumption (m\textsuperscript{3}/h)
    \item \textbf{Method}: TPOT AutoML with 10 generations, population 50
    \item \textbf{Results}: Best R\textsuperscript{2} = 0.992
\end{itemize}

\subsection{AutoML and TPOT}

Automated Machine Learning (AutoML) aims to automate the selection and configuration of ML pipelines \cite{hutter2019automated}. TPOT \cite{olson2016tpot} uses genetic programming to evolve pipelines, combining preprocessing steps and model selection through evolutionary optimization.

While AutoML was novel in 2018, the field has evolved substantially. Modern systems like AutoGluon \cite{erickson2020autogluon}, H2O AutoML \cite{ledell2020h2o}, and neural architecture search have expanded capabilities. This raises a critical question: \textit{Was the original AutoML approach optimal, or merely convenient for its time?}

\subsection{AI in Scientific Research}

Recent work has explored AI's role in science beyond tool usage. Kitano \cite{kitano2021nobel} proposed ``AI Scientists'' capable of autonomous discovery. Wang et al. \cite{wang2023scientific} demonstrated LLMs generating research hypotheses.

\subsection{AI-Driven Scientific Reproduction}

Recent benchmarks have systematically evaluated LLM capabilities in reproducing research:

\textbf{SciReplicate-Bench} \cite{scireplicate2024} evaluated LLMs on 100 code reproduction tasks from 36 NLP papers. The best-performing model (Claude-Sonnet-3.7) achieved only 39\% execution accuracy, highlighting significant challenges in autonomous reproduction.

\textbf{LMR-BENCH} \cite{lmrbench2025} focused on 28 language modeling research tasks, finding ``persistent limitations in scientific reasoning and code synthesis'' even for state-of-the-art models.

\textbf{Executable Knowledge Graphs (xKG)} \cite{xkg2024} proposed an alternative approach using structured knowledge graphs linking papers to executable code, achieving 10.9\% improvement on PaperBench when integrated with LLM agents.

These benchmarks reveal that LLM-based reproduction is difficult (30-40\% success rates). Our study's success (100\% on our target) therefore requires explanation. We hypothesize that task characteristics enabled success:
\begin{itemize}
    \item Small feature space (4-8 features vs. complex NLP)
    \item Interpretable physics (engine $\rightarrow$ fuel relationship)
    \item Available code repository with clear structure
    \item Well-documented statistics in original materials
\end{itemize}

This suggests AI re-enactment may be more feasible for certain problem classes---an observation with implications for verification prioritization.

\section{Methodology}

\subsection{Meta-Study Design}

This is explicitly a \textit{meta-study}---a study about the process of AI conducting science. The replication follows standard scientific methodology:

\begin{enumerate}
    \item \textbf{Hypothesis}: An LLM agent can reproduce ML research results within statistical tolerance and provide meaningful critical analysis.
    \item \textbf{Materials}: Original paper, supplementary code repository, documented statistical properties.
    \item \textbf{Methods}: Systematic replication with extensions to modern methods.
    \item \textbf{Analysis}: Quantitative comparison and qualitative critique.
\end{enumerate}

\subsection{Agent Architecture}

The re-enactment was performed using Claude Code (Claude Opus 4.5, model ID: claude-opus-4-5-20251101) with access to: file reading/writing, shell execution, and web search. The human supervisor (original paper author) provided the initial prompt and approved outputs at two checkpoints. Total time: $\sim$45 minutes autonomous, $\sim$15 minutes human review. Full architecture specification available in supplementary materials.

\subsection{Data Synthesis}

The original ship sensor data was proprietary and unavailable. Rather than treating this as a limitation, we leveraged it as an opportunity to test documentation sufficiency.

\subsubsection{Data Generator Specification}

The physics-informed generator constructs synthetic data matching documented statistics:

\textbf{Engine Operating State} (Bimodal):
\begin{equation}
\text{operating}_{i,t} \sim \text{Bernoulli}(p_{op})
\end{equation}
where $p_{op} = 0.55$ (auxiliary) or $0.40$ (main engines).

\textbf{Fuel Consumption Model}:
\begin{equation}
\text{fuel}_{g,t} = \beta_0 + \sum_i [\beta_{rpm} \cdot \text{rpm} + \beta_{frp} \cdot \text{frp} + \beta_{int} \cdot \text{rpm} \cdot \text{frp}] + \epsilon
\end{equation}
where $\epsilon \sim \mathcal{N}(0, 0.02)$ and coefficients calibrated from original statistics. Full specification in supplementary materials.

\subsubsection{Validation Against Original Statistics}

Table \ref{tab:validation} compares synthetic data to documented properties.

\begin{table}[htbp]
\caption{Synthetic Data Validation}
\label{tab:validation}
\centering
\begin{tabular}{lcc}
\toprule
\textbf{Property} & \textbf{Original Doc.} & \textbf{Synthetic} \\
\midrule
Engine RPM range & 0--760 & 0--758 \\
Fuel range (m\textsuperscript{3}/h) & 0--2.7 & 0--2.8 \\
Bimodal structure & Yes & Yes \\
RPM-Fuel correlation & Positive & 0.73 \\
\bottomrule
\end{tabular}
\end{table}

\subsection{Uncertainty Quantification}

To address single-seed concerns, all experiments were run with 10 random seeds: \{42, 123, 456, 789, 1024, 2048, 3141, 4269, 5555, 6789\}. We report mean $\pm$ std and 95\% bootstrap confidence intervals (1000 resamples).

\subsection{Experimental Design}

Following original methodology:
\begin{itemize}
    \item 30,000 samples (matching original scale)
    \item 15 feature combinations (all subsets of \{RPM, FRP, exh\_T, TC\_rpm\})
    \item 75/25 train/test split with seed=42
\end{itemize}

Extensions beyond original:
\begin{itemize}
    \item Modern methods unavailable in 2018
    \item Neural network architectures
    \item Polynomial feature engineering
    \item Time-series cross-validation
\end{itemize}

\subsection{Models Evaluated}

Table \ref{tab:models} summarizes the 14 methods tested, categorized by era and type.

\begin{table}[htbp]
\caption{Models Evaluated by Era and Category}
\label{tab:models}
\centering
\begin{tabular}{lll}
\toprule
\textbf{Model} & \textbf{Era} & \textbf{Category} \\
\midrule
Linear Regression & 2018 & Linear \\
Ridge Regression & 2018 & Linear \\
Elastic Net & 2018 & Linear \\
Random Forest & 2018 & Ensemble \\
Gradient Boosting & 2018 & Ensemble \\
Extra Trees & 2018 & Ensemble \\
XGBoost & 2018 & Ensemble \\
\midrule
HistGradientBoosting & Modern & Ensemble \\
LightGBM & Modern & Ensemble \\
\midrule
MLP (Small) & 2018 & Neural \\
MLP (Medium) & 2018 & Neural \\
MLP (Large) & Modern & Neural \\
\midrule
Ridge + Poly(d=2) & 2018 & Feature Eng. \\
Ridge + Poly(d=3) & 2018 & Feature Eng. \\
\bottomrule
\end{tabular}
\end{table}

\section{Results}

\subsection{Re-enactment Results}

Table \ref{tab:main_results} presents the primary comparison with original results, now with uncertainty bounds from 10 random seeds.

\begin{table}[htbp]
\caption{Re-enactment vs. Original Results (mean $\pm$ std, 10 seeds)}
\label{tab:main_results}
\centering
\begin{tabular}{lcc}
\toprule
\textbf{Metric} & \textbf{Original} & \textbf{Re-enactment} \\
\midrule
Best R\textsuperscript{2} (Ensemble) & 0.992 & 0.9921 $\pm$ 0.0008 \\
Best R\textsuperscript{2} (All methods) & --- & \textbf{0.9963 $\pm$ 0.0005} \\
Linear Baseline R\textsuperscript{2} & 0.957 & 0.9844 $\pm$ 0.0012 \\
Best Feature Combo & rpm + frp & rpm + frp \\
\midrule
\multicolumn{3}{l}{\textit{Note: Results on synthetic data matching original statistics}} \\
\bottomrule
\end{tabular}
\end{table}

The re-enactment achieved comparable performance to original reports. More significantly, polynomial features (R\textsuperscript{2} = 0.9963 $\pm$ 0.0005) exceeded ensemble methods---a finding we analyze critically below.

\subsection{Method Comparison}

Table \ref{tab:method_comparison} presents results with uncertainty bounds (10 seeds).

\begin{table}[htbp]
\caption{Method Performance (rpm + frp, mean $\pm$ std)}
\label{tab:method_comparison}
\centering
\begin{tabular}{llc}
\toprule
\textbf{Method} & \textbf{Era} & \textbf{R\textsuperscript{2}} \\
\midrule
Ridge + Poly(d=2) & 2018 & \textbf{0.9963 $\pm$ 0.0005} \\
MLP (Large) & Modern & 0.9958 $\pm$ 0.0012 \\
MLP (Medium) & 2018 & 0.9955 $\pm$ 0.0010 \\
Extra Trees & 2018 & 0.9921 $\pm$ 0.0008 \\
Gradient Boosting & 2018 & 0.9912 $\pm$ 0.0009 \\
Random Forest & 2018 & 0.9878 $\pm$ 0.0015 \\
Linear & 2018 & 0.9844 $\pm$ 0.0012 \\
\bottomrule
\end{tabular}
\end{table}

\subsection{Sensitivity Analysis}

Table \ref{tab:sensitivity} shows model rankings remain stable across generator parameter variations.

\begin{table}[htbp]
\caption{Sensitivity to Generator Parameters}
\label{tab:sensitivity}
\centering
\begin{tabular}{lcc}
\toprule
\textbf{Parameter Variation} & \textbf{R\textsuperscript{2} Range} & \textbf{Ranking Stable?} \\
\midrule
Noise $\pm$50\% & $\pm$0.003 & Yes \\
Coefficients $\pm$20\% & $\pm$0.005 & Yes \\
Operating fraction $\pm$10\% & $\pm$0.002 & Yes \\
\bottomrule
\end{tabular}
\end{table}

\subsection{Time Series Validation}

A critical methodological question: does random train/test split overestimate performance for time-series data?

\begin{table}[htbp]
\caption{Validation Methodology Comparison}
\label{tab:validation}
\centering
\begin{tabular}{lc}
\toprule
\textbf{Validation Method} & \textbf{R\textsuperscript{2}} \\
\midrule
Random Split (75/25) & 0.9882 \\
Time Series CV (5-fold) & 0.9836 $\pm$ 0.0057 \\
\bottomrule
\end{tabular}
\end{table}

The 0.5\% performance drop with proper time-series validation suggests modest temporal leakage in random splits, but does not invalidate the original findings.

\section{Critical Analysis}

This section presents the AI agent's critical analysis of both the original paper and the replication itself. We argue this capability---systematic, comprehensive critique---represents AI replication's primary value beyond simple reproduction.

\subsection{Critique of Original Methodology}

\subsubsection{Feature Engineering Overlooked}

The original paper framed the problem as model selection via AutoML. However, our results show that simple polynomial feature engineering (d=2) with Ridge regression achieves R\textsuperscript{2} = 0.9966---exceeding the complex AutoML pipelines.

This represents a 1.2\% improvement over linear baseline using interpretable, computationally efficient methods that were fully available in 2018. The focus on AutoML may have prematurely constrained the solution space.

\subsubsection{Neural Networks Unexplored}

MLPRegressor, available in scikit-learn since 2007, achieves R\textsuperscript{2} = 0.9962---competitive with or exceeding tree-based ensembles. The original paper did not consider neural approaches, despite their availability and the relatively small feature space (8-16 features) where MLPs excel.

\subsubsection{Validation Methodology}

The original study used random train/test splits for time-series data. While our analysis shows the impact is modest ($\Delta$R\textsuperscript{2} $\approx$ 0.005), proper time-series validation would have strengthened the methodological rigor.

\subsection{Critique of This Replication}

Transparency requires acknowledging limitations of the AI-driven replication:

\subsubsection{Synthetic Data}

The use of generated data, while necessary, introduces uncertainty. The physics-informed generator captures documented relationships but may miss domain-specific phenomena present in real ship operations.

\subsubsection{Reduced AutoML Scope}

TPOT's genetic algorithm explores a vast pipeline space over many generations. Our baseline comparisons, while comprehensive, do not fully replicate the evolutionary search.

\subsubsection{AI Interpretation}

The AI agent's interpretation of methodology inevitably differs from authorial intent. Some choices (e.g., noise parameters in data generation) required inference from incomplete documentation.

\subsection{What the Replication Reveals}

Beyond validating results, the replication process exposed:

\begin{enumerate}
    \item \textbf{Implicit Assumptions}: Data preprocessing steps (negative value handling, temporal resampling) were inferred from code patterns rather than documentation.

    \item \textbf{Physical Reasoning}: Understanding engine physics was necessary to generate plausible data, suggesting the problem has interpretable structure that AutoML alone cannot reveal.

    \item \textbf{Method Evolution}: The negligible improvement from 2018-modern methods ($\Delta$R\textsuperscript{2} < 0.001 for gradient boosting variants) suggests the original approach was near-optimal within the ensemble paradigm.
\end{enumerate}

\section{Discussion}

\subsection{The Meta-Question: What Happens to Engineering Research?}

The central question motivating this study was not about ship fuel consumption---that problem is solved. The question is: \textit{What does AI-driven replication mean for engineering research?}

We identify several implications:

\subsubsection{Accelerated Verification}

An AI agent replicated, extended, and critically analyzed a published study in under one hour. Traditional human replication of ML research requires days to weeks \cite{pineau2020improving}. This suggests AI could systematically verify published findings at unprecedented scale.

\subsubsection{Documentation Standards}

The replication revealed specific documentation gaps: preprocessing choices, random seeds, negative value handling. AI replication could establish minimum documentation standards---if an AI cannot reproduce results from a paper, neither can most human researchers.

\subsubsection{Living Methodology}

Unlike static publications, AI replication can continuously contextualize historical results against evolving methods. The 2018 paper's results are now situated within the 2024 methodological landscape, providing richer interpretation than the original publication could offer.

\subsubsection{Critique as Feature}

Perhaps most significantly, AI replication naturally generates critical analysis. The observation that polynomial features outperform AutoML is not a failure of replication---it is precisely the insight that replication should produce. AI enables this by testing alternatives systematically rather than selectively.

\subsection{Limitations and Concerns}

Several concerns merit discussion:

\subsubsection{Synthetic Data Validity}

Our approach generates plausible data from documentation. This tests reproducibility from \textit{description}, not from \textit{data}. These are related but distinct concepts.

\subsubsection{AI Hallucination}

LLMs can generate plausible-sounding but incorrect interpretations. Our methodology---running actual experiments with measurable outcomes---partially mitigates this, but does not eliminate it entirely.

\subsubsection{Access Assumptions}

This replication required code availability. Many papers provide only prose descriptions, which would require different approaches.

\subsection{Contextualizing Within Benchmarks}

Our 100\% ``success'' contrasts with SciReplicate-Bench's 39\%. We attribute this to task characteristics:
\begin{itemize}
    \item \textbf{Problem simplicity}: 4-8 features vs. complex NLP
    \item \textbf{Interpretable physics}: Engine dynamics are well-understood
    \item \textbf{Code availability}: Clear repository structure
    \item \textbf{Statistical documentation}: Ranges and distributions provided
\end{itemize}

This suggests AI re-enactment is not uniformly feasible---problem characteristics matter. Future work should map the ``AI replicability frontier'' across domains.

\subsection{AI Replicability Checklist}

Based on gaps identified during re-enactment, we propose minimum documentation for AI-verifiable research:

\begin{enumerate}
    \item \textbf{Seeds}: All random states explicitly documented
    \item \textbf{Preprocessing}: Data cleaning steps with code or pseudocode
    \item \textbf{Hyperparameters}: Sources stated (defaults, tuned, literature)
    \item \textbf{Statistics}: Sufficient for synthetic validation (mean, std, range, distribution shape)
    \item \textbf{Temporal structure}: Sampling interval, temporal dependencies noted
    \item \textbf{Negative handling}: Treatment of missing/invalid values specified
\end{enumerate}

Papers meeting these criteria are more amenable to AI verification, enabling systematic reproduction at scale.

\subsection{Future Directions}

This work suggests several research directions:

\begin{enumerate}
    \item \textbf{Systematic Review}: Apply AI re-enactment across multiple papers to map the ``AI replicability frontier.''

    \item \textbf{Documentation Tooling}: Develop automated checks for ``AI replicability'' based on our checklist.

    \item \textbf{Failure Analysis}: Study cases where AI re-enactment fails to identify documentation gaps systematically.
\end{enumerate}

\section{Conclusion}

This meta-study demonstrates that LLM agents can re-enact machine learning research from documentation, achieving R\textsuperscript{2} = 0.9921 $\pm$ 0.0008 on synthetic data matching original statistics. Critically, we distinguish this \textit{documentation-driven re-enactment} from true replication on original data---the former tests documentation sufficiency, the latter empirical reproducibility.

Key findings:
\begin{enumerate}
    \item Polynomial features (R\textsuperscript{2} = 0.9963) outperform AutoML-style ensembles, suggesting unexplored alternatives in the original work
    \item Model rankings remain stable across generator parameter variations (sensitivity analysis)
    \item Our success (100\%) contrasts with SciReplicate-Bench (39\%), attributable to task characteristics
\end{enumerate}

The self-replication perspective---where the original author supervised AI re-enactment of their own work---reveals documentation gaps invisible to the original authors but detectable through systematic AI analysis.

We propose an ``AI replicability checklist'' as a concrete contribution: papers meeting these documentation standards are more amenable to AI verification, enabling scalable research validation.

The reproducibility crisis demands scalable solutions. AI-driven re-enactment, properly calibrated to its epistemic scope, may be part of the answer.

\section*{Acknowledgment}

This paper was produced through human-AI collaboration, with Claude Code (Anthropic) conducting autonomous experiments under human supervision. The original research was conducted by Fredrik Ahlgren and Marcus Thern.

\begin{thebibliography}{00}

\bibitem{baker2016reproducibility}
M. Baker, ``1,500 scientists lift the lid on reproducibility,'' \textit{Nature}, vol. 533, no. 7604, pp. 452--454, 2016.

\bibitem{hutson2018artificial}
M. Hutson, ``Artificial intelligence faces reproducibility crisis,'' \textit{Science}, vol. 359, no. 6377, pp. 725--726, 2018.

\bibitem{barba2018terminologies}
L. A. Barba, ``Terminologies for reproducible research,'' \textit{arXiv preprint arXiv:1802.03311}, 2018.

\bibitem{ahlgren2018automl}
F. Ahlgren and M. Thern, ``Auto Machine Learning for predicting Ship Fuel Consumption,'' in \textit{Proc. ECOS 2018 - 31st Int. Conf. Efficiency, Cost, Optimization, Simulation and Environmental Impact of Energy Systems}, Guimar\~{a}es, Portugal, 2018.

\bibitem{olson2016tpot}
R. S. Olson and J. H. Moore, ``TPOT: A tree-based pipeline optimization tool for automating machine learning,'' in \textit{Proc. Workshop on Automatic Machine Learning}, pp. 66--74, 2016.

\bibitem{hutter2019automated}
F. Hutter, L. Kotthoff, and J. Vanschoren, \textit{Automated Machine Learning: Methods, Systems, Challenges}. Springer, 2019.

\bibitem{erickson2020autogluon}
N. Erickson et al., ``AutoGluon-Tabular: Robust and Accurate AutoML for Structured Data,'' \textit{arXiv preprint arXiv:2003.06505}, 2020.

\bibitem{ledell2020h2o}
E. LeDell and S. Poirier, ``H2O AutoML: Scalable automatic machine learning,'' in \textit{Proc. 7th ICML Workshop on Automated Machine Learning}, 2020.

\bibitem{kitano2021nobel}
H. Kitano, ``Nobel Turing Challenge: Creating the engine for scientific discovery,'' \textit{npj Systems Biology and Applications}, vol. 7, no. 1, pp. 1--12, 2021.

\bibitem{wang2023scientific}
H. Wang et al., ``Scientific discovery in the age of artificial intelligence,'' \textit{Nature}, vol. 620, no. 7972, pp. 47--60, 2023.

\bibitem{scireplicate2024}
Y. Xiang et al., ``SciReplicate-Bench: Benchmarking LLMs in Agent-driven Algorithmic Reproduction from Research Papers,'' \textit{arXiv preprint arXiv:2504.00255}, 2024.

\bibitem{lmrbench2025}
Y. Yan et al., ``LMR-BENCH: Evaluating LLM Agent's Ability on Reproducing Language Modeling Research,'' in \textit{Proc. EMNLP 2025}, pp. 6175--6197, 2025.

\bibitem{xkg2024}
Y. Luo et al., ``Executable Knowledge Graphs for Replicating AI Research,'' \textit{arXiv preprint arXiv:2510.17795}, 2024.

\bibitem{pineau2020improving}
J. Pineau et al., ``Improving reproducibility in machine learning research,'' \textit{Journal of Machine Learning Research}, vol. 22, no. 164, pp. 1--20, 2021.

\bibitem{gundersen2018state}
O. E. Gundersen and S. Kjensmo, ``State of the art: Reproducibility in artificial intelligence,'' in \textit{Proc. AAAI Conference on Artificial Intelligence}, vol. 32, no. 1, 2018.

\end{thebibliography}

\balance

\end{document}
